\section{Thí nghiệm hạn chót có tiền đăng ký (thăm dò)}
\label{muc:tang2}

Mục \ref{muc:haunghiem} và \ref{muc:hoito} chỉ là giải thích hậu nghiệm: quy tắc được dựng và
``kiểm'' trên cùng dữ liệu. Mục này thay chúng bằng một thí nghiệm thật: cài luật hạn chót vào bộ
điều phối, \emph{niêm phong} dự đoán theo từng ô trước khi chạy, rồi đối chiếu với kết cục, kể cả
trên một tập ô chưa từng dùng để dựng phương pháp. Mọi lần chạy có trôi dạt vẫn không qua bộ kiểm
chứng của dự án, nên kết quả mang nhãn \emph{thăm dò}.

\subsection{Thiết kế}

\textbf{Bốn họ, cùng trần 30 giây.} Cổng hành vi $C$ là cấu hình C1 như mọi mục trước
(\texttt{conduct-30}). Hạn chót một chiều $M^{+}$ (\texttt{dl-}$\ell$) bỏ trần ngân sách giờ trả và
thay bằng luật: loại mọi phương án làm giờ trả của một khách muộn hơn lời hứa \emph{đầu tiên} của
khách ấy quá $\ell$ giây, $\ell \in \{30, 60, 120\}$. Hạn chót hai chiều $M$ (\texttt{dl2-30}) loại
thêm phương án làm giờ trả sớm hơn lời hứa đầu quá 30 giây. Cổng kết quả $V$ (\texttt{visible-30})
là C1 với cơ sở tính khách nhìn thấy. Khóa pha và mọi chiều khác giữ nguyên như C1.

\textbf{Cài đặt.} Luật hạn chót mặc định tắt; khi tắt, bộ chạy cho kết quả giống từng byte với bộ
chạy cũ ở mọi chính sách không dùng bộ giải và không đổi quyết định thành công nào ở chính sách dùng
bộ giải (kiểm bằng 963 ca kiểm thử và hai vòng rà soát độc lập). Khác với mọi khóa hiện có, vốn so
phương án với mốc ngoại sinh $e$ nên trôi dạt không bao giờ làm vướng, hạn chót so với lời hứa đầu
$p_0$, nên riêng trôi dạt cũng làm vướng được --- kể cả trên phương án không làm gì. Bộ chạy mới cũng
trích lý do loại của chính phương án không làm gì khi một xe mất hết phương án.

\textbf{Hai panel.} Panel 2A là panel chính 20 ô, khung sáng 8:00--10:00 của năm ngày 14--18/11, ở
hai mức trôi dạt $k \in \{0, 1\}$. Panel 2B là panel phát triển 16 ô của hai ngày 12--13/11, chưa
từng dùng để dựng quy tắc dự đoán, gồm 8 ô khung sáng và 8 ô khung tối 17:00--19:00. Khung tối là
nơi Hệ quả \ref{hq:chieutroi} cho dự đoán phân biệt, vì hệ số thời gian đi đo thật \emph{giảm}
trong khung: ngày 12/11 hệ số tăng nhẹ lên $1{,}015$ lần ở phút thứ 15 rồi giảm dần còn $0{,}837$ lần
cuối khung.

\textbf{Biến kết cục.} Một ô là \emph{vô nghiệm do cam kết} khi lần chạy dừng với thông báo lỗi mà
lý do được trích là một luật cam kết --- ngân sách, khóa pha hay hạn chót --- đối chiếu đúng hằng
thông điệp trong mã nguồn. Ô dừng vì bất biến cửa sổ đón (mục \ref{muc:batbien}) trước khi vướng
hạn chót được xếp là \emph{không} vô nghiệm do cam kết, một cách xếp chống lại giả thuyết.

\subsection{Đối chứng, quỹ đạo tham chiếu và niêm phong}

\textbf{Đối chứng.} Trước mọi nhánh mới, bộ chạy mới được chạy lại hai cấu hình đã có ở mức $k = 0$
trên đủ 20 ô panel 2A. Ở 36 ô hoàn tất, chuỗi kế hoạch được gán, tập lời hứa công bố ở từng khung
quyết định, mã lý do và số khách được nhận trùng \emph{tuyệt đối} với lần chạy cũ. Cổng kết quả vô
nghiệm ở đúng bốn ô 16-r2, 16-r4, 17-r1, 17-r3 như bảng \ref{bang:doicho}, và tổng khách được nhận
trên 16 ô còn lại là $1.294$, đúng con số ở bảng \ref{bang:coso}. Vì bộ chạy mới trích lý do của
chính phương án không làm gì, lần này thấy được trực tiếp điều Định lý \ref{dl:ketqua} nói: ở cả bốn
ô, phương án không làm gì bị loại vì vượt ngân sách.

\textbf{Quỹ đạo tham chiếu.} Cổng hành vi chạy trước ở $k = 1$ trên cả hai panel để có quỹ đạo gần
dự phòng: panel 2A có 9 ô hoàn tất --- trùng từng quyết định với lần chạy cũ --- và 11 ô dừng vì bất
biến cửa sổ đón; panel 2B có 13 ô hoàn tất và 3 ô khung sáng dừng vì bất biến ấy. Không ô nào vô
nghiệm do cam kết, đúng phần (a) của Định lý \ref{dl:hocong}. Một lần chạy thử ở khung tối xác nhận
chín hệ số áp trong mô phỏng trùng từng giá trị với tệp đo thật.

\textbf{Dự đoán theo ô.} Trên quỹ đạo tham chiếu cùng ô, cùng mức trôi dạt, với mỗi khách tính cực
đại chạy $M^{+}$, $M$ và tổng chạy $V$ của chuỗi giờ trả đã công bố; $\tau$ là thời điểm đầu tiên một
khách vượt ngưỡng, $T$ là khung quyết định cuối nếu tham chiếu dừng. Quy tắc: $\tau \le T$ thì đoán
vô nghiệm do cam kết; không vượt và tham chiếu hoàn tất thì đoán không; không vượt trước khi tham
chiếu dừng thì đoán dừng vì bất biến cửa sổ đón (dự đoán yếu). Đây là một \emph{heuristic}, không
phải hệ quả của định lý: nhánh hạn chót đi quỹ đạo khác tham chiếu ngay sau lần cắt đầu tiên.

\textbf{Niêm phong.} Tiền đăng ký, 200 dự đoán và mã chạy, phân loại, phân tích được băm SHA-256 lúc
06:03:52 UTC ngày 24/9/2026, khi chưa có một lần chạy nào của các nhánh được dự đoán. Một lỗi được
phát hiện \emph{trước} niêm phong và được báo cáo: bản dự đoán đầu tiên đoán nhầm mọi ô tham chiếu
hoàn tất là vô nghiệm, do phép so $\infty \le \infty$; bản ấy được giữ lại, không niêm phong và không
dùng. Tiền đăng ký ghi trước ngưỡng xác nhận, điều kiện phủ định và một bảng dự đoán hướng cho khung
tối lấy từ Hệ quả \ref{hq:chieutroi}.
